\chapter{数据、标签与标注工程}
\label{chap:edge-ai-data-label-engineering}

视觉项目最昂贵的错误通常不是少训练几个 epoch，而是数据无法代表目标现场、标签语义含糊或切分泄漏。模型会忠实学习数据中最容易利用的相关性，包括背景、设备水印、拍摄 session 和人工标注习惯。本章把数据视为可版本化的产品输入，而不是一个不断覆盖的图片目录。

\section{从目标域而不是公开榜单出发}

数据设计首先描述部署分布：相机、镜头、安装角度、工作距离、分辨率、ISP、曝光、光照、温度、运动、背景、对象版本和操作者行为。公开数据可提供通用预训练和任务启动，但不能代替目标域证据。ImageNet 与 COCO 等数据集推动了通用分类和检测研究\cite{imagenet,coco}，产品仍需验证它们的类别、视角和许可是否适合当前目标。

推荐把数据来源分为四层：

\begin{enumerate}
  \item \textbf{通用预训练数据}：学习纹理、形状和通用对象特征；
  \item \textbf{领域公开数据}：接近目标行业、视角或任务；
  \item \textbf{真实目标域数据}：来自计划部署的设备、位置和流程；
  \item \textbf{固定金标回归集}：独立审核、永不混入训练，用于发布比较。
\end{enumerate}

一个来源样本数量很大，不等于覆盖充分。连续视频的十万帧可能只代表一个设备、一个背景和几分钟光照。

\section{公开数据集的发现、获取与组合}

公开数据集应被当作版本化外部依赖，而不是从搜索结果下载后直接合并。搜索时同时检查数据索引、论文/benchmark 页面和维护者的一手下载页；聚合站或镜像可以用于发现和容灾，但不能替代原始 provenance。候选选择至少比较：

\begin{itemize}
  \item 类别、任务和标签粒度是否真正匹配，还是只能间接推断；
  \item sensor、分辨率、视角、对象尺度、环境和困难条件是否接近目标域；
  \item 是否提供 subject/sequence/site/device 等可用于安全分组的身份；
  \item 标注规范、漏标/噪声、重复、固定 benchmark split 和已知排除规则；
  \item 当前一手访问地址、版本、注册/鉴权、归档大小、checksum 和下载可恢复性；
  \item 许可证、用途限制、署名、再分发和隐私条件是否允许计划用途。
\end{itemize}

为每个版本建立 acquisition manifest，状态从 discovered、verified-page、sampled、downloaded、checksum-verified、extracted 到 audited 单向推进。最小字段可以是：

\begin{lstlisting}[caption={公开数据集 acquisition manifest 字段示意}]
dataset_id: source-name-v1
source_page: https://first-party.example/dataset
version: v1
checked_at: 2026-08-02
archives:
  - {name: images.zip, sha256: "verified-digest"}
access: {method: https, registration: false}
roots: {raw: data/public/source-name/v1/raw,
        derived: data/public/source-name/v1/derived}
license: {id: source-declared-license, reviewed: true}
status: checksum-verified
conversion_revision: git:0123456789ab
\end{lstlisting}

示例中的地址和摘要只是字段结构，正式清单必须保存当次核验的一手 URL、实际摘要、访问日期和手工步骤，凭据不能写进 manifest 或 shell 历史。原始归档只读保存，解压和格式转换进入派生目录；下载重试不得覆盖已经校验的版本。

混合前为每个来源建立显式映射：\texttt{source class/field} 到 canonical class/field、坐标/单位变换、valid/missing/ignore、审核需求和转换 revision。宽泛类别不能未经证据拆成细分类，图像级标签也不能冒充对象框。数据卡应记录来源、映射、支持量、用途和已知缺口\cite{datasheets-for-datasets}。

跨来源先做完全重复、感知近重复和 group 检查，再分配内部 split。保留官方 benchmark test 的原始用途，不能因为拿到了标签或镜像就把它混入训练。大型预训练模型可能已经见过公开 benchmark，因此产品晋级仍以独立 reviewed 目标域集合为准。最终组合应按来源记录采样权重和指标，防止巨大通用数据淹没稀缺目标域证据或产生不可见的负迁移。

\section{采集覆盖矩阵}

在采集前建立场景维度，并选择需要覆盖的组合：

\begin{table}[H]
  \centering
  \caption{视觉数据采集覆盖矩阵示例}
  \begin{tabular}{p{24mm}p{45mm}p{43mm}p{25mm}}
    \hline
    维度 & 典型取值 & 高风险缺口 & 记录字段 \\
    \hline
    设备与光学 & sensor、镜头、焦距、ISP 版本 & 只用一台样机 & \texttt{device\_id}、校准版本 \\
    环境 & 白天、夜间、逆光、频闪、温湿度 & 实验室均匀照明 & site、照度、曝光 \\
    对象 & 型号、颜色、磨损、批次、尺寸 & 训练与量产批次不同 & object/batch ID \\
    姿态与尺度 & 距离、旋转、遮挡、边缘位置 & 只采中心正面 & ROI、姿态、可见比例 \\
    负样本 & 相似件、手、工具、反光、空场景 & 只有容易正样本 & negative reason \\
    时间 & 班次、季节、设备老化和维护后 & 一次性短采集 & session、timestamp \\
    \hline
  \end{tabular}
\end{table}

覆盖矩阵不是要求枚举所有笛卡尔积，而是让风险显式化。对失败代价高的切片应主动采集，对低风险组合可以抽样。采集应用应保存原始媒体和元数据，避免只留下已经 resize、压缩或裁剪后的训练输入。

\section{采集合同与小规模试点}

大规模采集前先冻结 capture/sensor contract。对于图像或视频，至少记录：

\begin{itemize}
  \item sensor、镜头、安装和硬件 revision，固件、驱动、ISP 与采集程序版本；
  \item 原始分辨率、像素格式、bit depth、颜色空间/range、编码和帧率；
  \item 曝光、gain、白平衡、焦点是固定还是自动，以及自动策略状态；
  \item 单调时钟/墙上时钟来源、时间精度、丢帧和重启后的 sequence 规则；
  \item 相机内参、畸变、标定残差和适用分辨率；多传感器还要记录外参与同步误差；
  \item site/device/session/object/lot 等稳定 ID，不能只依赖可变文件名；
  \item 隐私、授权、操作安全、停止条件和数据保留策略。
\end{itemize}

改变镜头、焦点、分辨率、crop/binning 或关键 ISP 变换后，旧标定不能默认继续有效。产品声明要泛化到新设备、站点或批次时，这些对象同时也是覆盖统计和 group split 的独立单位。

正式扩量前，用少量 session 完成一次端到端 pilot：采集、断点恢复、传输、checksum、解码、元数据 join、标注平台导入、样本读取、一次训练冒烟和归档。Pilot 需要证明：

\begin{enumerate}
  \item 标签在目标图像中确实可观察，而不是依赖采集人员的场外知识；
  \item 媒体颜色、方向、时间戳和标定与合同一致；
  \item 采集应用能在掉电、存储满或网络中断后安全恢复；
  \item 单 session 的时间、容量、人工成本和失败率可估算；
  \item 计划覆盖与实际覆盖能够在采集过程中反馈给操作员。
\end{enumerate}

无法通过小规模 pilot 的缺口，扩大到十万张图后只会更昂贵。后处理无法恢复没有记录的曝光、设备身份或同步信息，也不能让不可观察的标签重新出现。

\section{不可变原始数据与 Manifest}

原始数据区只追加，不覆盖、不原地重编码。清洗、裁剪、标注和增强都通过 manifest 指向来源。每个样本至少包含：

\begin{itemize}
  \item 稳定 sample ID、相对路径、内容摘要、宽高和编码；
  \item device/site/subject/session/video/frame 等分组键；
  \item 标签、逐字段 valid、审核状态和标签来源；
  \item 原始采集条件、转换链和父样本 ID；
  \item 数据许可、访问级别、保留期限和敏感性分类。
\end{itemize}

JSONL 适合教学项目，因为每行可独立解析和 diff：

\begin{lstlisting}[caption={多任务视觉样本的简化 JSONL 结构}]
{"sample_id":"site01-s03-f00142",
 "image":"raw/site01/session03/frame00142.jpg",
 "width":1920,"height":1080,
 "group":{"site":"site01","session":"s03","device":"cam02"},
 "labels":{"scene":"wrong_pose","scene_valid":true,
           "objects":[{"class":"connector","bbox_xywh":[412,280,330,210],
                       "bbox_valid":true,"pose":"reversed","pose_valid":true}]},
 "label_source":"human_reviewed","review_state":"accepted"}
\end{lstlisting}

正式数据版本应包含样本清单摘要、schema 版本、split 规则、生成代码 revision 和统计报告。复制一个目录并命名为 \texttt{final2} 无法重建实验。

\section{入库、隔离与质量门}

采集完成不等于数据已经可发布。推荐的数据流为：

\begin{center}
\texttt{capture $\rightarrow$ staging $\rightarrow$ integrity/quality audit}
\\
\texttt{$\rightarrow$ immutable raw version $\rightarrow$ derived/annotation views}
\end{center}

Staging 阶段检查预期文件、大小、checksum、媒体解码、分辨率/格式、帧数、时间戳单调性、元数据 join 和 session 完整性。损坏、缺元数据、配置错误、标定不匹配或策略不允许的数据进入 quarantine，并保留原因；不能静默删除后仍把采集成功率报告为 100\%。

原始数据版本发布前还应完成：

\begin{itemize}
  \item 随机样本、最低质量尾部、稀有场景和覆盖缺口的 contact sheet/视频抽查；
  \item 完全重复、近重复和重复场景聚类，保存代表样本映射；
  \item blur、过暗、饱和、遮挡、压缩损坏等质量统计，但不把统一阈值当作所有设备真理；
  \item 标定残差、同步误差和设备/对象/session 分组完整性；
  \item 原始媒体、抽帧、裁剪、脱敏和训练输入使用不同目录，并由 parent sample ID 关联。
\end{itemize}

先去重和确定 group，再分配 split。若先随机切分后去重，相似副本可能已经泄漏到测试集，而且后续删除会改变各集合分布。

\section{标签 Schema 先于标注界面}

标签定义要回答对象、属性、坐标、单位、互斥关系和无效状态。以检测框为例，需要明确：

\begin{itemize}
  \item 坐标是像素、归一化还是物理单位；
  \item 使用 \texttt{xyxy}、\texttt{xywh} 或中心点格式；
  \item 边界是闭区间还是半开区间，是否允许超出图像；
  \item 遮挡对象标可见框还是完整估计框；
  \item 截断、模糊、太小和无法判断如何表示；
  \item 属性属于对象、图像、轨迹还是采集 session。
\end{itemize}

同义词必须映射到唯一类别 ID。类别顺序一旦进入模型 ABI，就不能仅修改显示名称而不提升 schema/ABI 版本。

\section{缺失、不可判断与真实负类}

以下三种状态必须分开：

\begin{description}
  \item[真实负类] 标注者已检查该字段，并确认目标不存在或属性为否；
  \item[不可判断] 图像证据不足，例如严重遮挡、过曝或分辨率不够；
  \item[尚未标注] 当前任务没有审核该字段，不能用于正负监督。
\end{description}

把“尚未标注”默认为负类，会让扩展新任务时已有数据持续产生错误梯度。把“不可判断”强行分到某一类别，则会把人为猜测固化为模型边界。

\section{标注规范与一致性}

标注规范应包含正例、反例、边界案例和优先级规则。对难以文字描述的类别，建立代表性图册。标注流程至少包括：

\begin{enumerate}
  \item 小规模试标，收集分歧并修订 schema；
  \item 两人独立标注一部分重叠样本；
  \item 统计分类一致率、框 IoU、关键点距离或 mask 一致性；
  \item 由裁决者解决系统性歧义，并把决定写回规范；
  \item 正式标注过程中持续抽检，而不是结束时一次验收。
\end{enumerate}

高一致率也可能来自两位标注者共同误解。抽检应包含原始场景上下文和产品专家审核，尤其是高风险类别。

\section{分组切分与泄漏}

训练、验证、测试必须按与泛化声明一致的 group 切分。若产品要泛化到新设备，就按设备隔离；要泛化到新用户，就按用户隔离；视频任务通常至少按视频或 session 隔离。

随机帧切分会让相邻帧、相同背景和同一对象同时出现在训练与测试，得到虚高指标。重复与近重复检查可组合：

\begin{itemize}
  \item 文件摘要发现完全重复；
  \item 感知哈希发现缩放、压缩和轻微裁剪副本；
  \item 图像 Embedding 发现语义近重复；
  \item 元数据检查同一 subject/session/device 是否跨 split；
  \item 人工检查最高相似对，确认阈值不会误删真正困难样本。
\end{itemize}

测试集在项目早期冻结。反复根据测试结果选择模型、阈值和数据方向，会把测试集逐渐变成验证集；此时需要新的独立 holdout。

\section{类别不平衡与支持量}

长尾数据需要同时报告样本数、对象数、主体数和 session 数。某类别有一万张相邻帧但只来自两个对象，其独立支持仍然很弱。

常见处理路径按优先级为：

\begin{enumerate}
  \item 重新采集关键少数类和强负样本；
  \item 以 group 和类别控制采样，避免同一 session 支配 batch；
  \item 使用适当的数据增强；
  \item 再考虑 class weight、focal loss 或重采样；
  \item 用每类指标和校准结果确认处理是否真的改善，而非只降低训练 Loss。
\end{enumerate}

损失权重不能创造不存在的视觉变化，也不能补偿错误标签。

\section{数据增强必须保持标签语义}

增强用于模拟部署中合理变化。每项增强都要回答“现场是否可能发生”和“标签如何同步变化”。

\begin{itemize}
  \item 水平翻转会改变左右关键点顺序、文字方向和某些不对称产品姿态；
  \item 随机 crop 可能截断目标，需要更新框、mask、关键点和 truncation；
  \item 旋转会改变角度符号、坐标和 padding；
  \item 颜色增强应尊重目标 ISP、曝光和照明范围，不能制造物理上不可能的颜色；
  \item mosaic/mixup 等强增强可能改善检测泛化，也可能破坏细小缺陷和图像级状态语义。
\end{itemize}

训练、验证和部署的确定性预处理必须来自同一契约；随机增强只存在于训练分支。

\section{真实、公开与合成数据如何混合}

合成数据适合补充稀有、危险、昂贵或需要精确几何标签的场景，但它是待验证的数据源，不是“免费金标”。根据任务选择最简单且能保持证据的方法：

\begin{itemize}
  \item 2D compositing 控制对象、背景、尺度和遮挡；
  \item 3D/物理仿真生成 pose、depth、mask、flow、track 和多传感器真值；
  \item 参数化渲染生成 OCR、屏幕、仪表和版面；
  \item 带 edit mask 的生成式编辑扩展外观或缺陷，但几何关键标签需独立验证；
  \item domain randomization 在有边界的真实范围和压力尾部内改变光照、材质、相机和环境\cite{domain-randomization}。
\end{itemize}

生成器必须版本化资产、场景、sensor profile、随机分布、相关约束、seed、后端和标签映射。防止渲染边框、水印、固定背景、资产 ID 或不可能组合成为捷径。标签经过 crop、畸变、resize 和遮挡后要与图像执行完全相同的几何变换。

合成数据的晋级实验至少比较：

\begin{enumerate}
  \item real-only baseline；
  \item synthetic-only 诊断，判断模型学到何种模拟特征；
  \item real + synthetic，在相同训练预算或明确额外计算下比较；
  \item 不同混合比例、先合成预训练再真实微调、source-balanced batch 等策略。
\end{enumerate}

最终判断只看独立 reviewed 真实目标域集合。合成样本不能进入真实 test，也不能因为视觉效果逼真就宣称提升。若某种生成模式没有改善目标门槛，或增加关键切片 FP，应停止使用并保留最小决策记录。

\section{预标注、教师模型与主动学习}

当数据量增大时，人工应从“从零绘制”转向“审核模型候选”。不同教师可以负责不同任务：开放词汇检测发现新类别，领域检测器提供主要框，分割/跟踪模型传播视频标签，视觉语言模型提供保守的图像级候选。

启用自动化前先建立一个小型 reviewed gold seed，覆盖关键类别、空场景、歧义、困难条件和主要设备域。每个预标注 provider 使用稳定接口返回 sample/object identity、任务、几何/文本/属性、原始分数、校准置信度、模型/配置版本、prompt/ROI 和失败原因。未检出必须与“确认负类”分开。

CVAT、Label Studio、FiftyOne 或自研平台只是审核和展示适配器，canonical schema/manifest 才是事实源。平台导入导出必须按稳定 sample ID 和 schema version 合并，支持幂等重试、原子保存、冲突处理、审核状态和人工修改历史。不能按文件名猜身份，也不能让下一次 provider 导入覆盖已经 reviewed 的字段。

教师输出必须记录模型版本、原始分数、支持任务和人工修改。人工删除的错误候选要留下 rejection，避免下一轮重复添加。

多 provider 融合必须是可审计的 proposal 生成过程，而不是把多个分数简单平均后写回金标。先按任务适用的类别、几何、对象身份和时间约束匹配候选；只把真正支持该任务的 provider 计入共识分母，同一 provider 的类别别名、TTA 或重复框要先去重。不同 provider 的原始 score 不在同一概率空间，必须分别按类别和目标域校准，不能直接比较 0.9 与 0.9。

融合结果应保留每个贡献者、匹配关系、未匹配候选和分歧原因。高一致候选可以进入快速审核队列，低一致或高业务风险候选进入强制复核；没有任何 provider 检出的区域仍不能自动成为强负类。人工已审核字段具有最高优先级，融合器只能追加新 proposal，不能静默覆盖、删除或降级人工结论。

视频标注先对相邻帧去重或聚类，选择多样 anchor frame，审核 anchor 后使用 tracker、光流或分割传播。身份切换、突变、低置信区间和片段边界进入人工队列，不能把传播结果无条件扩展到整段视频。

主动学习队列不应只选择最低置信样本，否则会被大量相似模糊帧淹没。可组合：

\begin{equation}
  P=U+D+R+N+E-\lambda C,
\end{equation}

其中 $U$ 为不确定性，$D$ 为教师分歧，$R$ 为稀有性，$N$ 为域新颖度，$E$ 为当前模型错误或高损失，$C$ 为重复惩罚。权重应在固定金标集上验证，不把公式当作通用常数。

自动接受伪标签前，需在独立 reviewed 集上按类别校准阈值，并保留固定比例抽检。伪标签始终标记为伪标签，不能改名为人工金标，也不能进入验证集。

标注自动化是否有效，需要与纯人工基线在相同质量门下比较：每人工小时审核的样本/对象/字段数、accept/edit/reject、漏标率、随机抽检错误、类别/域覆盖和固定回归集的增益。高置信 proposal 很多但人工修改率没有下降，不能宣称节省了标注成本。

\section{数据质量报告}

每个数据版本至少输出：

\begin{itemize}
  \item split、类别、任务 valid、group、设备、站点和困难切片统计；
  \item 缺图、损坏、越界框、非法类别、空标注和重复检查；
  \item 标签来源、审核状态、修改率和抽检错误率；
  \item 数据转换、压缩、resize 和坐标往返测试；
  \item 已知覆盖缺口、许可证和访问限制。
\end{itemize}

统计脚本应在 CI 或数据发布流程中重复执行。只在 notebook 中展示一次图表，无法阻止后续数据版本重新引入泄漏或非法标签。

\section{实践：构建第一个可审计数据版本}

围绕贯穿项目完成以下步骤：

\begin{enumerate}
  \item 定义四分类或检测任务的 schema、valid 和 unknown；
  \item 对至少一个公开数据源建立 acquisition manifest、标签映射和接受/拒绝理由；
  \item 从至少三次独立采集 session 获取数据，保留原始图像和元数据；
  \item 编写 manifest 生成与校验脚本；
  \item 按 session 分组生成 train/validation/test；
  \item 运行完全重复、感知哈希和 group 泄漏检查；
  \item 双人复核一小部分样本并记录分歧；
  \item 在 reviewed gold seed 上校准至少一个预标注 provider，记录逐类 Precision/Recall、accept/edit/reject 和人工小时；
  \item 输出数据卡：用途、来源、覆盖、限制、许可和不建议用途；
  \item 冻结 \texttt{v001}，后续修正通过新版本完成。
\end{enumerate}

通过标准是任何成员都能从 manifest 重建相同 split，并解释每个样本为什么属于该集合、标签由谁确认、哪些目标域仍未覆盖。

实践至少保留以下可审计产物，而不是只留下一个训练目录：

\begin{itemize}
  \item 版本化的 capture contract、label schema 与字段迁移说明；
  \item 原始数据 checksum 清单、采集 session 索引和 quarantine 记录；
  \item 逐样本 manifest、确定性 split 配置与生成命令；
  \item 机器可读质量报告、人工抽检记录和已知覆盖缺口；
  \item 数据卡、许可证/授权边界，以及从原始版本到派生视图的 provenance；
  \item 能在干净环境中重放的校验入口及其预期退出码。
\end{itemize}

最后主动做一次数据流水线故障演练：把同一 group 放入训练集和测试集、篡改一个媒体文件、注入未知类别、删除必填元数据，并确认 Gate 分别报告泄漏、checksum 失败、schema 失败和完整性失败。每个失败都应阻止数据版本晋级，同时把样本送入可追踪的修复或隔离路径。只有正常数据可通过、异常数据会稳定失败，才算完成这个实践。
